% ═══════════════════════════════════════════════════════════════
% KURZTEST: Reaktion von Metallen mit Sauerstoff
% Thema 2.4 | Chemie Klasse 8 | 15 Minuten | 20 BE
% ═══════════════════════════════════════════════════════════════

\documentclass[11pt, a4paper]{article}

\usepackage[utf8]{inputenc}
\usepackage[T1]{fontenc}
\usepackage[ngerman]{babel}

\usepackage{helvet}
\renewcommand{\familydefault}{\sfdefault}

\usepackage[margin=2cm, top=2.5cm, bottom=2cm]{geometry}
\usepackage{enumitem}
\usepackage{array}
\usepackage{tabularx}
\usepackage{booktabs}

\usepackage{amsmath, amssymb}

\usepackage{fancyhdr}
\usepackage{lastpage}

\usepackage{xcolor}
\usepackage{tcolorbox}
\tcbuselibrary{skins}

% ═══════════════════════════════════════════════════════════════
% FARBEN
% ═══════════════════════════════════════════════════════════════

\definecolor{chemie}{HTML}{2a7a4b}
\definecolor{hellgrau}{HTML}{F5F5F5}
\definecolor{mittelgrau}{HTML}{888888}
\definecolor{dunkelgrau}{HTML}{444444}

% ═══════════════════════════════════════════════════════════════
% KOPFZEILE
% ═══════════════════════════════════════════════════════════════

\pagestyle{fancy}
\fancyhf{}
\renewcommand{\headrulewidth}{0.5pt}
\renewcommand{\footrulewidth}{0pt}
\fancyhead[L]{\textbf{Chemie} | Kurztest: Metalloxide}
\fancyhead[R]{Klasse 8 | \_\_\_\_\_\_\_\_\_\_\_\_\_\_\_\_\_\_\_\_}
\fancyfoot[R]{\small\textcolor{mittelgrau}{\thepage}}

% ═══════════════════════════════════════════════════════════════
% BEFEHLE
% ═══════════════════════════════════════════════════════════════

\newcommand{\punkte}[1]{\hfill\textcolor{mittelgrau}{/#1}}
\newcommand{\luecke}[1]{\underline{\hspace{#1}}}

\setlength{\parindent}{0pt}
\setlength{\parskip}{6pt}

% ═══════════════════════════════════════════════════════════════
% DOKUMENT
% ═══════════════════════════════════════════════════════════════

\begin{document}

% ═══════════════════════════════════════════════════════════════
% KOPFBEREICH
% ═══════════════════════════════════════════════════════════════

\begin{tcolorbox}[colback=hellgrau, colframe=hellgrau, boxrule=0pt, arc=0pt, left=8pt, right=8pt, top=6pt, bottom=6pt]
\begin{tabular}{@{}p{8cm}p{6cm}@{}}
\textbf{\large Kurztest: Reaktion von Metallen mit Sauerstoff} & \textbf{Name:} \luecke{4cm} \\[0.3em]
Thema 2.4 | 15 Minuten | 20 BE & \textbf{Datum:} \luecke{3cm}
\end{tabular}
\end{tcolorbox}

\vspace{0.8em}

% ═══════════════════════════════════════════════════════════════
% AUFGABE 1: Begriffe (AFB I)
% ═══════════════════════════════════════════════════════════════

\textbf{Aufgabe 1:} \textit{Nenne} die passenden Fachbegriffe. \punkte{4}

\begin{enumerate}[label=\alph*), leftmargin=2em, itemsep=0.5em]
    \item Die Reaktion eines Stoffes mit Sauerstoff heißt \luecke{4cm}.
    \item Eine chemische Verbindung aus einem Element und Sauerstoff heißt \luecke{3cm}.
    \item Die langsame Oxidation von Metallen im Alltag nennt man \luecke{3cm}.
    \item Die schützende Oxidschicht bei Aluminium nennt man \luecke{3cm}.
\end{enumerate}

\vspace{1em}

% ═══════════════════════════════════════════════════════════════
% AUFGABE 2: Wortgleichung (AFB I)
% ═══════════════════════════════════════════════════════════════

\textbf{Aufgabe 2:} \textit{Ergänze} die Wortgleichungen. \punkte{3}

\begin{enumerate}[label=\alph*), leftmargin=2em, itemsep=0.8em]
    \item Magnesium + Sauerstoff $\longrightarrow$ \luecke{4cm}
    \item \luecke{3cm} + Sauerstoff $\longrightarrow$ Eisenoxid
    \item Kupfer + \luecke{3cm} $\longrightarrow$ Kupferoxid
\end{enumerate}

\vspace{1em}

% ═══════════════════════════════════════════════════════════════
% AUFGABE 3: Massenerhaltung (AFB II)
% ═══════════════════════════════════════════════════════════════

\textbf{Aufgabe 3:} \textit{Berechne} und \textit{erkläre}. \punkte{5}

12 g Magnesium reagieren mit 8 g Sauerstoff.

\begin{enumerate}[label=\alph*), leftmargin=2em, itemsep=0.3em]
    \item Wie viel Gramm Magnesiumoxid entstehen? \punkte{1}

    \vspace{0.3em}
    Antwort: \luecke{2cm} g

    \item Erkläre, warum das Produkt schwerer ist als das Metall allein. \punkte{4}

    \vspace{2.5cm}
\end{enumerate}

% ═══════════════════════════════════════════════════════════════
% AUFGABE 4: Korrosion (AFB II)
% ═══════════════════════════════════════════════════════════════

\textbf{Aufgabe 4:} \textit{Vergleiche} die Oxidation von Eisen und Aluminium. \punkte{4}

\begin{tabular}{|p{3cm}|p{5.5cm}|p{5.5cm}|}
\hline
& \textbf{Eisen} & \textbf{Aluminium} \\
\hline
Oxidschicht-Struktur & & \\[1.2em]
\hline
Schutzwirkung? & & \\[1.2em]
\hline
\end{tabular}

\vspace{1em}

% ═══════════════════════════════════════════════════════════════
% AUFGABE 5: Formelgleichung (AFB II) ★★
% ═══════════════════════════════════════════════════════════════

\textbf{Aufgabe 5 (★★):} \textit{Schreibe} die Formelgleichung für die Verbrennung von Magnesium. \punkte{4}

\vspace{0.3em}
\hspace{2em} \luecke{1.5cm} Mg + \luecke{1.5cm} $\longrightarrow$ \luecke{1.5cm} \luecke{2cm}

\vspace{0.5em}
{\small Hinweis: Sauerstoff kommt als O$_2$-Molekül vor. Gleiche die Gleichung aus.}

\vspace{1.5em}

% ═══════════════════════════════════════════════════════════════
% BEWERTUNG
% ═══════════════════════════════════════════════════════════════

\begin{tcolorbox}[colback=hellgrau, colframe=hellgrau, boxrule=0pt, arc=0pt, left=8pt, right=8pt, top=4pt, bottom=4pt]
\small
\begin{tabular}{@{}p{10cm}p{4cm}@{}}
\textbf{Bewertung:} & \textbf{Erreicht:} \_\_\_\_ / 20 BE \\[0.3em]
20--18: sehr gut | 17--14: gut | 13--10: befriedigend | 9--6: ausreichend & \textbf{Note:} \_\_\_\_\_\_\_\_
\end{tabular}
\end{tcolorbox}

\end{document}
